Ein schwimmender Körper taucht so tief ein, dass die Auftriebskraft seine Gewichtskraft gerade aufhebt:
\begin{align}
 F_\text{A} &= F_\text{G} \quad\Rightarrow\quad \rho_\text{W}\,V\,g = m\,g
\end{align}
Er verdrängt also gerade seine eigene Masse an Wasser ($\rho_\text{W} = \rhoWO$, FoTaBe S.~187):
\begin{align}
 V &= \VSF \\
   &= \frac{\mS}{\rhoW} \\
   &= \VS \approx \result{\VSlP{-1}{2}}
\end{align}
Dass die Schüssel aus Aluminium ist, spielt keine Rolle: Massgebend ist nur ihre Masse.
